\section{Introduction}
\label{sec:intro}
Real-world query workloads increasingly rely on expressive natural language (NL) interfaces that seamlessly blend predicates over structured relational attributes (e.g., a movie's release year, duration, or genre) with unstructured textual attributes (e.g., user reviews, synopses, or customer feedback). Evaluating unstructured predicates within such hybrid queries requires nuanced semantic interpretation. While traditional text processing and classical sentiment analysis techniques~\cite{Chandan2025ACS} have long been used for text classification, they fundamentally fail to handle ambiguous, indirect, or misleading language, nor can they reliably perform accurate identification of complex opinions expressed through sarcasm, negation, or implicit tone. Consequently, incorporating Large Language Models (LLMs) as relational semantic query operators has emerged as a necessity for capturing complex semantic intent. However, because LLM inference incurs high computational latency and monetary cost, there is an urgent need for efficient query processing methods that can evaluate natural language queries over combined structured and unstructured predicates without sacrificing execution speed or answer quality.

To illustrate such queries, consider the following representative question over a relational movie database:
\begin{quote}
\textit{``Which movies released in $1998$ received user reviews containing sarcastic, backhanded, or contextually negative feedback about the film's execution?''}

\smallskip
\noindent\textbf{Example Review:} \textit{``If your idea of a thrilling Friday night is sitting in a pitch-black theater for  hours watching paint dry while listening to a self-important monologue on existential dread, then this absolute masterpiece of pure boredom is definitely the movie for you.''}
\end{quote}
Answering this query requires evaluating a structured predicate $\sigma_S$ over schema attributes ($\text{year} = 1998$) alongside an unstructured semantic predicate $\sigma_U$ over free-text reviews. Executing such hybrid queries presents three major challenges. First, traditional relational query engines (SQL) cannot express or execute $\sigma_U$, as there is no static vocabulary to match against;  opinions manifest in diverse, open-ended ways that evade simple keyword search. Second, even if data can be preprocessed, there is no fixed relational schema or index that can organize reviews by overall tone or sentiment. Third, delegating $\sigma_U$ to an LLM operator introduces severe performance bottlenecks: naively invoking an LLM for every candidate tuple or attempting unguided semantic joins over raw collections leads to overwhelming wall-clock latency and monetary expense, making full-table LLM scans intractable for interactive systems.

Existing frameworks offer partial solutions to this problem, but fall at opposite ends of a the cost-quality trade-off. SUQL~\cite{liu2023suql} adopts a pushdown filtering strategy. It compiles the structured predicate $\sigma_S$ into a standard SQL filter and evaluates the semantic predicate $\sigma_U$ by issuing one row-wise LLM call ($\mathtt{answer()}$) per candidate row surviving $\sigma_S$. While this yields high precision because every candidate receives full model context, LLM invocations scale linearly $O(|R_\sigma|)$ with surviving rows, causing severe latency delays when structured filters are moderately selective. Trummer~\cite{trummer2025semanticjoin} proposes a batched block-join approach that avoids structured pushdown entirely. Instead, it partitions input data into blocks and prompts the LLM to execute both $\sigma_S$ and $\sigma_U$ simultaneously over a batch of candidate rows in a single call. Although batching drastically reduces the  number of LLM calls, it causes a significant drop in accuracy and recall because the model is overburdened with performing joins, structured filtering, and sentiment analysis in a single forward pass.
Other semantic execution engines, such as ELEET~\cite{urban2024eleet} (multi-modal operator fusion and cost-based operator reordering) and Stretto~\cite{sanmartino2025stretto} (Bayesian threshold calibration and KV-cache execution ladders), further highlight the need for principled semantic operator optimization.

To overcome the limitations of these existing partial solutions, we introduce a \textbf{cascaded semantic join architecture} that combines early relational pushdown with multi-level LLM execution. We propose two cascaded framework variants, \textsc{CasSuql} and \textsc{CasTrummer}, built upon two foundational contributions:
\begin{enumerate}
    \item \textbf{Structured Filter Pushdown and Batched Execution:} We enforce strict early execution of structured filters ($\sigma_S$) to eliminate non-matching candidates for free before any LLM prompt is constructed. Furthermore, \textsc{CasTrummer} batches surviving rows into multi-item calls, amortizing prompt prefix overhead across candidate items.
    \item \textbf{Calibrated Multi-Stage Cascading:} We layer a lightweight, low-cost cheap pre-filtering stage before the full-context, larger LLM operator. To decide whether the cheap filter's score is trustworthy, we employ a Bayesian calibration model using a Beta posterior over empirical recall and precision bounds ($\hat{r}_\alpha(\tau), \hat{p}_\alpha(\tau)$) fit on a minimal calibration sample (20 rows, stratified by cheap score, whose expensive labels are kept as final answers). This defines a mathematically sound ambiguity interval $[\tau_-, \tau_+]$: candidates outside the band are resolved immediately by the cheap filter, while only ambiguous units are escalated to expensive, full-context LLM verification. A side of the band that the sample cannot certify stays open, so the cascade falls back to the expensive model instead of guessing.
\end{enumerate}

Empirical evaluations across the \textbf{IMDb} and \textbf{Amazon Fashion} benchmarks demonstrate the benefit of this cascaded approach. Specifically, cascading the batched join baseline (\textsc{CasTrummer}) achieves a superior cost-quality trade-off: on \textbf{IMDb}, it cuts LLM calls by $81\%$ relative to \textsc{SUQL} at about the same wall-clock time, while delivering the highest $F_1$ score ($0.713$, against $0.501$ for \textsc{SUQL} and $0.340$ for \textsc{Trummer}) and the fewest calls among all evaluated approaches; on \textbf{Amazon Fashion}, it is the cheapest method in calls, time and dollars ($83\%$ fewer calls and $20\%$ less time than \textsc{SUQL}) and improves the $F_1$ of \textsc{Trummer} by $45\%$. \textsc{CasSuql} preserves the quality of \textsc{SUQL} ($F_1$ $0.500$ against $0.501$ on \textbf{IMDb}, $0.796$ against $0.779$ on \textbf{Amazon Fashion}) and never issues more expensive calls than \textsc{SUQL}, but it is slower ($1.8$--$1.9\times$), because the cheap model is not faster per call than the expensive one. By combining structured pushdown, batched verification and a certified cheap stage, \textsc{CasTrummer} lies on the cost--$F_1$ Pareto front of both datasets.

We summarize the related work in \S\ref{sec:related}. We introduce our query model, baseline architectures and our proposed solution in \S\ref{sec:preliminaries}. Our proposed cascaded solutions are described in \S\ref{sec:contributions}. The results of our empirical evaluation are reported in \S\ref{sec:eval}. We conclude and discuss future work in \S\ref{sec:conclusion}.
